\section{The Harvest Vault}
\label{sec:lp}

The Harvest Vault gives passive LPs pool-style UX while the operator runs an active lending strategy, bounded by a publicly auditable mandate and verified per-epoch by zero-knowledge proof.

\subsection{Operator and Mandate}

The operator's strategy is bounded by a \textbf{mandate}, a policy object committed on-chain as a hash at vault deployment. The mandate specifies:
\begin{itemize}
  \item \textbf{Allowed collateral set:} asset addresses the operator may accept as loan collateral.
  \item \textbf{Per-collateral LTV ceiling:} maximum credit multiplier the operator may offer on each collateral type.
  \item \textbf{Per-tenor rate floor:} minimum annualized rate the operator must charge, ensuring the vault never underwrites below a specified spread.
  \item \textbf{Portfolio concentration caps:} maximum exposure per borrower, per collateral type, and per market as a fraction of vault NAV.
  \item \textbf{Liquidation reserve requirement:} minimum idle capital maintained at all times for liquidation response and withdrawal settlement.
\end{itemize}

The mandate constrains what the operator \emph{can} do, not what they \emph{must} do; within its bounds, the operator retains full discretion. The full mandate specification is public and auditable.

\subsection{Compliance Proof}

Each epoch, the operator produces a zero-knowledge proof attesting that every action taken satisfied the mandate. The proof takes prior vault state, the epoch's operator actions, and the mandate hash as public inputs, and outputs a \textbf{compliance fact} registered in the facts registry (Section~\ref{sec:settlement}):
\begin{equation}
  \begin{split}
    f_{\mathit{compliance}} &= \khash(h_{\mathit{prog}},\, \mathit{chainId},\, \mathit{vault}, \\
    &\qquad \mathit{epoch}_{\mathit{prev}},\, \mathit{epoch}_{\mathit{next}},\, h_{\mathit{mandate}})
  \end{split}
  \label{eq:compliancefact}
\end{equation}
where $h_{\mathit{prog}}$ is the compliance program hash, $\mathit{vault}$ is the vault contract address, and $h_{\mathit{mandate}}$ is the committed mandate hash.

\subsection{Fee-Gating}

The operator earns a spread fee escrowed during each epoch. A valid compliance proof releases the escrow to the operator; otherwise it flows back into NAV. An invalid or missing proof does not halt loan settlement or affect borrowers; only the operator's cut is at risk.

\subsection{NAV and Shares}

LPs deposit USDC and receive ERC-20 shares. NAV comprises idle USDC, debt-token claim values, escrowed hedging buffers, and accrued interest, minus USDC ring-fenced for pending withdrawals. The share price reflects NAV per outstanding supply, with a virtual offset against inflation attacks. Appreciation comes from fixed-rate interest on vault-funded loans, buffer surplus releases on externally-sourced loans, and liquidation proceeds. Impairments reduce NAV directly; LPs absorb residual risk through share-price adjustment.

\subsection{Epoch Withdrawal}
\label{subsec:withdrawal}

Instant withdrawal is incompatible with capital deployed in active positions, per-loan buffer escrows, and transient liquidation bridges. The vault queues exits in a four-step flow:

\begin{enumerate}
  \item \textbf{Request.} The LP burns vault shares and receives \emph{withdrawal receipt} tokens for the current epoch.
  \item \textbf{Epoch elapses.} After the epoch duration, anyone can advance the epoch.
  \item \textbf{Settlement.} A permissionless call snapshots the share price for all pending requests in that epoch and ring-fences the corresponding USDC, removing it from the vault's deployable idle balance.
  \item \textbf{Redemption.} The LP burns receipt tokens and receives the ring-fenced USDC, net of a flat redemption fee.
\end{enumerate}

Until settlement, the LP remains exposed to pool gains and losses. The epoch delay prevents withdrawing at a stale NAV while a loss event is about to land.

\subsection{Revenue Model}

The portfolio blend between vault-funded and externally-sourced loans (Figure~\ref{fig:capital}) is jointly determined by the operator's mandate and the risk engine's pricing.

\begin{table}[ht]
  \centering
  \small
  \begin{tabular}{@{}llll@{}}
    \toprule
    \textbf{Party} & \textbf{Risk borne} & \textbf{Why} & \textbf{Compensation} \\
    \midrule
    Borrower       & Collateral decline     & Seeks leverage      & Fixed-rate USDC \\
    Ext.\ venue    & Smart-contract, util.  & Money market         & Variable rate \\
    Rate buffer    & Variable ${>}$ fixed   & Rate transform       & Surplus \\
    Option buffer  & Swap cost ${>}$ escrow & Execution risk       & Escrow ${\sim}$95\% \\
    Vault LP       & Default, exhaustion    & Absorbs residual     & Coupon or surplus \\
    Operator       & Mandate violation      & Active mgmt          & Fee, gated on proof \\
    \bottomrule
  \end{tabular}
  \caption{Risk decomposition. Each risk channel is routed to the participant with the comparative advantage to bear it. The rate buffer and option buffer rows apply only to externally-sourced loans.}
  \label{tab:risk}
\end{table}
